\section{Experimental Results and Performance Analysis}
This section presents the empirical evaluation of the proposed lightweight DRM system compared against the traditional full-stream encryption baseline. The benchmarks were executed across a heterogeneous test dataset using the automated orchestration architecture.

\subsection{Experimental Setup and Throughput Formulation}
The performance profiling was conducted on an Intel I7-12700H host system running Python 3.13. To mathematically evaluate processing efficiency at the player side, the operational Data Decryption Throughput ($T$) is formally defined as:
\begin{equation}
T = \frac{S_{\text{bytes}} \times 8}{t_{\text{decryption}} \times 10^6} \quad \text{(Mbps)}
\end{equation}
where $S_{\text{bytes}}$ represents the exact file payload scale in bytes, and $t_{\text{decryption}}$ denotes the total clock execution duration in seconds measured via precision performance counters.

\begin{table*}[t]
\caption{Comprehensive Experimental Metrics: Baseline (Full) vs. Proposed (Selective)}
\label{tab:metrics}
\begin{center}
\begin{tabular}{lrrrrrrrr}
\toprule
\textbf{Video Asset Name} & \textbf{Size (MB)} & \textbf{Enc. Base (s)} & \textbf{Enc. Prop (s)} & \textbf{Dec. Base (s)} & \textbf{Dec. Prop (s)} & \textbf{CPU Base (\%)} & \textbf{CPU Prop (\%)} \\
\midrule
Trades in Space (orig)   & 1119.91 & 0.75442 & 2.44902 & 0.78234 & 0.05441 & 22.0 & 17.1 \\
Trades in Space (large)  & 104.50  & 0.08926 & 0.18721 & 0.06297 & 0.01272 & 23.9 & 10.5 \\
Trades in Space (medium) & 45.24   & 0.02099 & 0.07382 & 0.03587 & 0.01508 & 17.1 & 1.2  \\
Trades in Space (small)  & 27.28   & 0.01254 & 0.04321 & 0.02115 & 0.00912 & 14.5 & 1.1  \\
\bottomrule
\end{tabular}
\end{center}
\end{table*}

\begin{table}[t]
\caption{Decryption Throughput Comparison}
\label{tab:throughput}
\begin{center}
\setlength{\tabcolsep}{3pt}
\begin{tabular}{lrrr}
\toprule
\textbf{Asset Scale} & \textbf{Throughput Base} & \textbf{Throughput Prop} & \textbf{Gain Factor} \\
\midrule
Original (1.1 GB) & 12008.1 Mbps & 172669.1 Mbps & $14.38\times$ \\
Large (104 MB)    & 13922.4 Mbps & 68856.6 Mbps  & $4.94\times$  \\
Medium (45 MB)    & 10579.9 Mbps & 25122.3 Mbps  & $2.37\times$  \\
\bottomrule
\end{tabular}
\end{center}
\end{table}

\subsection{Cryptographic Latency and Computational Overhead}
The empirical results compiled during data execution show that the proposed design significantly outperforms the full-stream approach. Table I provides the data metrics collected during the comparative analysis.

\begin{figure}[htbp]
\centering
\includegraphics[width=0.48\textwidth]{decryption_latency.png}
\caption{Client Player Decryption Latency vs Video Asset Scale.}
\label{fig:dec_latency}
\end{figure}

\begin{figure}[htbp]
\centering
\includegraphics[width=0.48\textwidth]{throughput_rate.png}
\caption{Data Decryption Throughput Rate Performance Curves.}
\label{fig:throughput_plot}
\end{figure}

As plotted in Fig. 1 and Fig. 2, the decryption footprint at the client player shows a massive reduction in processing time. For the 1119.91 MB original video payload, the traditional framework required 0.78234 seconds to decipher the continuous block chains, whereas the proposed selective decoding loop completed execution in just 0.05441 seconds. This yields up to a 93.0\% reduction in client-side latency. The performance improvement occurs because the client demuxer completely bypasses non-structural P and B slices.

However, an engineering trade-off is observed within the server packaging stage. The encryption phase of the proposed architecture requires more time than the baseline (e.g., 2.44902 seconds vs. 0.75442 seconds for the original 1.1 GB file). This execution overhead is attributed to the sequential I/O cost of parsing individual byte arrays to isolate the Annex-B start codes. While the baseline full encryption processes the block arrays globally in a single optimized memory pass, the selective packager must inspect the bitstream sequentially to validate individual NAL headers via bitwise masking operations. 

Nevertheless, because video packaging is entirely an \textit{offline ingestion task} performed only once during asset onboarding, this upfront server-side overhead is highly acceptable. It represents a favorable trade-off for eliminating runtime hardware rendering bottlenecks and achieving significant performance gains on constraint client terminals during active playback.

\subsection{Visual Quality and Degradation Analysis}
To evaluate perceptual quality bounds under a simulated intercept vector, an unauthorized playback attempt was forced without fetching the key payload from the remote architecture.

\begin{table}[htbp]
\caption{Visual Quality Metrics Under Different Access Conditions}
\label{tab:security_metrics}
\begin{center}
\begin{tabular}{lrr}
\toprule
\textbf{Execution Scenario} & \textbf{PSNR Index (dB)} & \textbf{SSIM Value} \\
\midrule
Authorized (Valid CEK)    & $\infty$ (Perfect)       & 1.0000 \\
Unauthorized (No CEK)     & 8.4300 (Degraded)        & 0.1767 \\
\bottomrule
\end{tabular}
\end{center}
\end{table}

\begin{figure}[htbp]
\centering
\includegraphics[width=0.48\textwidth]{visual_degradation.png}
\caption{Perceptual Quality Visual Degradation Profile under Unauthorized Attempt.}
\label{fig:visual_degradation_plot}
\end{figure}

When a valid CEK is configured, the system yields clean rendering maps ($\text{PSNR} = \infty$, $\text{SSIM} = 1.0$). Conversely, under unauthorized decoding conditions, the intentional mutation of parameter sets and anchor tracking streams causes a catastrophic rendering collapse. 

As plotted in Fig. 3 and summarized in Table III, the visual presentation metrics drop to a highly broken state ($\text{PSNR} = 8.43\text{ dB}$, $\text{SSIM} = 0.1767$). Because the structural baseline sets are missing, standard media players cannot initialize macroblock residue calculations, rendering the output video as an unwatchable mix of static noise and corrupted macroblocks. This demonstrates sufficient practical protection for digital video distribution networks with minimal cryptographic processing cost.